\begin{abstract}
Phonetic speech transcription is crucial for fine-grained linguistic analysis and downstream speech applications. While Connectionist Temporal Classification (CTC) is a widely used approach for such tasks due to its efficiency, it often falls short in recognition performance, especially under unclear and nonfluent speech. In this work, we propose LCS-CTC, a two-stage framework for phoneme-level speech recognition that combines a similarity-aware local alignment algorithm with a constrained CTC training objective. By predicting fine-grained frame-phoneme cost matrices and applying a modified Longest Common Subsequence (LCS) algorithm, our method identifies high-confidence alignment zones which are used to constrain the CTC decoding path space, thereby reducing overfitting and improving generalization ability, which enables both robust recognition and text-free forced alignment. Experiments on both LibriSpeech and PPA demonstrate that LCS-CTC consistently outperforms vanilla CTC baselines, suggesting its potential to unify phoneme modeling across fluent and non-fluent speech. 
\end{abstract}

\begin{IEEEkeywords}
Phoneme, Transcription, Alignment, Clinical
\end{IEEEkeywords}